\documentclass[10pt]{article}
\usepackage[margin=2.0cm]{geometry}
\usepackage{graphicx}\usepackage{booktabs}\usepackage{amsmath}\usepackage{amssymb}
\usepackage{xcolor}\usepackage[hidelinks]{hyperref}\usepackage{enumitem}\usepackage{fontspec}
\setlength{\parskip}{4pt}\setlength{\parindent}{0pt}\setlist{itemsep=1pt,topsep=2pt}
\graphicspath{{../../output/08_two_color/cambaliza2000/pulses/delay_sweep/}}
\title{\vspace{-1.2cm}\textbf{Work to date: from the focus of few-cycle pulses\\ to two-color excitation with optical pulses}}
\author{Roland Albert A. Romero\\ \small National Institute of Physics, University of the Philippines Diliman}
\date{2026-11-11}
\begin{document}\maketitle\vspace{-0.6cm}

This document states what I have done so far, in the order in which it was done. Each part gives what was computed,
the main results against the paper it starts from, and where the details are. The numbers are those of the stored
results in \texttt{output/} and of the review documents in \texttt{docs/review/}.

\section*{Summary}
\begin{enumerate}
\item I reproduced the five figures of Romallosa, Bantang \& Saloma (2003) on the focus of few-cycle pulses with
the Richards--Wolf theory, and found the conventions that the paper leaves unstated.
\item I extended that focus to Gaussian input, annular pupils and an axicon, and found where the scalar theory stops
being accurate as NA increases.
\item I built a Monte Carlo model of the same focus, which gives the Debye--Wolf field without scattering, and
validated it through a scattering slab against six independent checks.
\item I reproduced Cambaliza \& Saloma (2000) and Lim \& Saloma (2002, 2003) on two-color excitation (2CE) with two
focused beams, and extended their results to high NA, Gaussian input, water and oil.
\item I extended 2CE to optical pulses. At 100 fs, pulsed 2CE is the same as with continuous beams. With facing
lenses and shorter pulses, 2CE forms where the two pulses meet, so its depth of focus is shorter.
\item I added a delay between the two pulses. A delay weakens 2CE through the temporal overlap of the pulses, and
with facing lenses at low NA it moves the 2CE along the axis instead.
\end{enumerate}

\section*{1. The focus of few-cycle pulses: Romallosa et al.\ (2003)}
Romallosa et al.\ compute the three-dimensional intensity near the focus of a 750-nm beam of Gaussian pulses, from
1 fs to continuous, with the Richards--Wolf vector theory, and compare it with the continuous beam through the
Linfoot criteria.
\begin{enumerate}
\item The five figures of the paper are reproduced. The structural content along the axis at 1 fs is 1.058, against
1.057 in the paper, and the transverse FWHM is 472.8 nm. The remaining differences are within 1\% for Fig.~4 and about
3\% for Fig.~3 at 1 fs.
\item The paper integrates the whole positive spectrum, from 0 to $2\omega_c$. The band it quotes, 483 nm to
1.67 $\mu$m, is the FWHM of the spectrum. With the FWHM band alone, 24\% of the spectral weight is lost and the
quoted values are not reproduced.
\item Three conventions had to be found from the figures: the panel labels of Fig.~2 are exchanged, Fig.~3 is an
integral along a line, and the Linfoot arrays are in optical coordinates, which is why the criteria do not depend on
NA.
\end{enumerate}
Details are in \texttt{docs/review/01\_romallosa/review\_01.pdf}.

\section*{2. Numerical aperture and Gaussian input}
\begin{enumerate}
\item At NA 0.1 the vector theory, the scalar Debye integral and the Airy pattern agree to $10^{-4}$.
\item As NA increases, the scalar theory departs from the vector theory first across the axis, along the
polarization. The FWHM differs by 1\% at NA 0.30 and by 5\% at NA 0.60 for linear polarization. Along the axis the
difference reaches 1\% at NA 1.15.
\item The difference is the longitudinal field $E_z$, which fills the zeros and widens the spot along the
polarization.
\item A Gaussian input with truncation parameter $\alpha = 4$ keeps the scalar theory accurate to a higher NA
(1\% at 0.40 and 5\% at 0.85), since it gives less weight to the marginal rays.
\end{enumerate}
Details are in \texttt{docs/review/02\_na\_apodization/review\_02.pdf}.

\section*{3. Annular pupils and the axicon}
\begin{enumerate}
\item The field of an annular pupil with obstruction ratio $\varepsilon$ is the field of the full pupil minus
that of the obstructed part. At low NA the depth of focus grows as $1/(1-\varepsilon^2)$, and the central lobe
narrows.
\item A thin annulus gives the $J_0^2$ profile of Durnin at low NA. At high NA the vector profile is narrower than
the scalar one (FWHM 2.06 against 2.25 in optical units).
\item The peak intensity falls as $(1-\varepsilon^2)^2$.
\item An axicon gives a focal segment with a Bessel profile whose scale depends on the position along the axis.
\end{enumerate}
Details are in \texttt{review\_03.pdf} and \texttt{review\_04.pdf}.

\section*{4. Few-cycle pulses with annular pupils}
\begin{enumerate}
\item The longer depth of focus of an annular pupil is kept with few-cycle pulses: the ratio of the depths of focus
at $\varepsilon = 0.99$ is 44.35 for a continuous beam and 44.33 at 1 fs.
\item The contrast of the transverse rings decreases with shorter pulses, as in Romallosa et al.\ for the circular
pupil.
\end{enumerate}
Details are in \texttt{review\_05.pdf}.

\section*{5. Monte Carlo model of the focus, with scattering}
\begin{enumerate}
\item Each photon is a plane-wave component of the converging field, and the photons are added with their phases.
Without scattering the sum is the Debye--Wolf integral: with $10^5$ photons the RMS difference from the
deterministic field is 0.0008.
\item The same holds for the Gaussian input, the annular pupil, the axicon and the vector field, each within 0.002.
\item Through a scattering slab, the unscattered part keeps the focus and loses amplitude, and the scattered part
forms a wide background. Six checks pass: the limit of no scattering, the Beer--Lambert law over the cone of rays,
energy, single scattering, the diffusion limit of a thick slab, and the diffusion solution for a point source with
absorption.
\item Each Monte Carlo result is the mean of at least ten independent runs, with the runs stored.
\end{enumerate}
Details are in \texttt{review\_06.pdf}. A comparison with an experiment or an independent full-wave calculation has
not been done.

\section*{6. Two-color excitation with continuous beams}
\subsection*{6.1 The papers}
In 2CE two beams of wavelengths $\lambda_1 = 750$ nm and $\lambda_2 = 656$ nm, with
$1/\lambda_e = 1/\lambda_1 + 1/\lambda_2$ and $\lambda_e = 350$ nm, are focused by two lenses whose axes make an
angle $\theta$. The fluorescence is $F_{2c} = \kappa I_1I_2$, and for two-photon excitation (2PE) at 700 nm it is
$F_{2p} = \kappa'|U_1+U_2|^4$ (Cambaliza \& Saloma, 2000). Lim \& Saloma (2002) displace the focus of the second
beam, and Lim \& Saloma (2003) give each beam a primary spherical aberration.

\subsection*{6.2 Replication}
\begin{enumerate}
\item Cambaliza \& Saloma: the profiles of Figs.~2 to 7 are reproduced at NA 0.5 for collinear ($\theta = 0$),
crossed ($\pi/2$) and facing ($\pi$) lenses, with circular and annular pupils. With crossed lenses the profiles
are the sharpest, and 2PE is narrower than 2CE with stronger side lobes, as in the paper. Several of the ratios
quoted in the text of the paper differ from the current work at $\theta = \pi/2$ (for example 0.67 against 0.49
for the 2PE spot at $\pi/2$ over that at 0), under the FWHM, the $1/e^2$ width and the first minimum, while the
widths plotted in Fig.~5 of the paper agree with the current work. The second beam is rotated in real space; the
rotation of the optical coordinates written in the papers does not reproduce the fits of Lim \& Saloma (2002).
\item Lim \& Saloma (2002): the peak 2CE fluorescence against the displacement of the focus follows the fits of the
paper. With crossed lenses and a displacement in the focal plane, the fit is $-0.063\,x_c^2 + 0.999$ against
$-0.06\,x_c^2 + 1$ in the paper. Along the axis it is $-0.045\,z_c^2 + 0.995$ against $-0.045\,z_c^2 + 0.986$.
The maps of fidelity, structural content and correlation quality behave as described in the paper.
\item Lim \& Saloma (2003): the peak is kept when the two aberration coefficients are equal at $\theta = 0$ and
opposite at $\theta = \pi$, and with crossed lenses the 2CE fidelity stays above 0.9 for coefficients up to 40, as
in the paper.
\end{enumerate}

\subsection*{6.3 Extension to high NA, Gaussian input and immersion}
\begin{enumerate}
\item As NA increases, the displacement of the foci that 2CE tolerates decreases as a power of NA: about
$\mathrm{NA}^{-1}$ across the axis and $\mathrm{NA}^{-2}$ along it.
\item The paraxial scalar model of the papers gives a larger tolerance along the axis than the vector model: by 7\%
at NA 0.5 and by 47\% near NA 0.95. The difference comes from the paraxial phase, not from the polarization.
\item The 2CE and 2PE profiles of the scalar model stay within 5\% of the vector ones up to NA 0.86 across the axis
and NA 0.43 along it, in air. In water the two limits are NA 1.14 and 0.57.
\item A Gaussian input ($\alpha = 4$) makes the tolerances about 1.4 times larger.
\item Crossed lenses of the same NA fit in air up to NA 0.71, where the two cones of rays touch.
\end{enumerate}
Paraxial scalar results are shown for NA up to 0.5 in all figures.

\section*{7. Two-color excitation with optical pulses}
\subsection*{7.1 Model}
Each beam is a transform-limited Gaussian pulse of intensity FWHM $\tau$, the sum of its frequencies focused by the
lens. The signal is the two-photon amplitude at the 350-nm level,
$A = A^{11} + 2A^{12} + A^{22}$ with $A^{ab} = \int E_aE_b\,e^{i\omega_et}dt$, and $F = |A|^2$. $2A^{12}$ takes one
photon from each beam (the 350-nm pairs, 750 + 656 nm). $A^{11}$ and $A^{22}$ take both photons from one beam (the
375-nm pairs, 750 + 750 nm, and the 328-nm pairs, 656 + 656 nm). For long pulses $F \to I_1I_2$, the result of
Cambaliza \& Saloma. The derivations and their checks are in
\texttt{docs/review/07\_two\_color\_pulses/review\_07.pdf}.

\subsection*{7.2 What was run}
The runs cover the three lens angles, NA 0.1 to 0.95 in air, to 1.25 in water and to 1.45 in oil, the scalar and
vector models, circular and annular pupils ($\varepsilon = 0.5, 0.9, 0.99$) and the Gaussian input
($\alpha = 4$). The pulse widths are 4 to 1000 fs without delay, and 5 to 100 fs with delays up to $2\tau$.
Unless stated, the numbers below are for the vector model in air with the circular pupil, both beams polarized
perpendicular to the plane of the lens axes. The table gives the continuous-beam widths that the pulsed results
are compared with.
\begin{center}\small\begin{tabular}{llccccc}\toprule
Continuous beams, FWHM ($\mu$m) & & NA 0.2 & 0.3 & 0.5 & 0.7 & 0.9\\\midrule
2CE, $\theta = 0$ and $\pi$ & across the axis & 1.29 & 0.86 & 0.51 & 0.36 & 0.27\\
 & along the axis & 22.0 & 9.67 & 3.33 & 1.56 & 0.81\\
 & along the axis, scalar model & 22.3 & 9.89 & 3.56 & 1.82 & 1.10\\
2CE, $\theta = \pi/2$ & along $x$ & 1.92 & 1.28 & 0.75 & 0.52 & \\
 & along $z$ & 1.68 & 1.12 & 0.66 & 0.46 & \\
2PE, $\theta = \pi/2$ & along $x$ and $z$ & 0.255 & 0.255 & 0.256 & 0.256 & \\
2PE, $\theta = \pi$ & along the axis, central lobe & 0.129 & 0.130 & 0.136 & 0.147 & 0.166\\\bottomrule
\end{tabular}\end{center}

\subsection*{7.3 Collinear lenses ($\theta = 0$)}
\begin{enumerate}
\item As NA increases, the 2CE spot and depth of focus decrease together, and 2CE and 2PE have nearly the same
profile. The depth of focus is 3.33 $\mu$m at NA 0.5 and 0.81 $\mu$m at NA 0.9.
\item The scalar model gives a longer depth of focus than the vector model, by 7\% at NA 0.5 and 36\% at NA 0.9.
Across the axis the two agree to about 1\% at NA 0.5.
\item With pulses, 2CE keeps the continuous-beam profile at all pulse widths computed, down to 4 fs, because the
two pulses travel together and overlap over the whole focus.
\item With a delay, the peak falls with the temporal overlap of the two pulses,
$\exp(-2\ln 2\,\Delta t^2/\tau^2)$: to one half at $\Delta t = 0.71\tau$ and one quarter at $\Delta t = \tau$,
at all NAs. The stored peaks follow this to $10^{-5}$ at 100 fs.
\item Below 20 fs, pulses with a fixed phase between the beams give fringes in the peak against the delay, with
periods of 2.19, 2.50 and 1.17 fs. Their swing is 0.92 of the peak at 5 fs and 0.04 at 20 fs, and it is the same
at all NAs, because with collinear lenses the pairs have no fringes in space.
\end{enumerate}

\subsection*{7.4 Crossed lenses ($\theta = \pi/2$)}
\begin{enumerate}
\item As NA increases, both 2CE widths decrease as about 1/NA, because the transverse spot of the other beam sets
each of them. The 2CE spot is 0.75 by 0.66 $\mu$m at NA 0.5.
\item The depth of focus is shorter than with collinear lenses, and this advantage decreases as NA increases: it
is 13 times shorter at NA 0.2 and 3.4 times at NA 0.7.
\item Two lenses of the same NA fit at right angles up to NA 0.71 in air and 0.94 in water, where their cones of
rays touch.
\item The 2PE central spot stays at 0.256 $\mu$m at all NAs, because the interference of the two beams sets it.
2PE is narrower than 2CE and has stronger side lobes, as in Cambaliza \& Saloma.
\item With pulses, the profiles match the continuous-beam ones from about 20 to 28 fs for 2CE. Below that the
depth of focus changes by a few percent, so crossed lenses keep the sharpest profile.
\item With a delay, the peak falls with the temporal overlap, to 0.250 to 0.253 at $\Delta t = \tau$.
\item Below 20 fs the fringes against the delay deepen with NA, because a smaller spot holds fewer of the fringes
in space: at 5 fs the swing along $x$ is 0.20 at NA 0.3 and 0.65 at NA 0.7.
\item With the polarization in the plane of the lens axes, the two fields are perpendicular at the focus and
the three kinds of pairs fall in different field components. The fringes are then about a quarter as deep
(amplitude about 0.05 against 0.21 at NA 0.5 and 5 fs), and the 2CE spot is wider (0.83 against 0.75 $\mu$m at
NA 0.5). The 350-nm pairs fall with the delay in the same way for both polarizations.
\end{enumerate}

\subsection*{7.5 Facing lenses ($\theta = \pi$)}
\begin{enumerate}
\item With continuous beams, 2CE has the same profile as with collinear lenses, since $I_1I_2$ does not depend
on the direction of beam 2. 2PE has a standing wave along the axis, with a central lobe of 0.13 to 0.17 $\mu$m.
\item With pulses, 2CE forms where the two pulses meet, a stretch of FWHM $G = c\tau/\sqrt2\,n$ (4.2 $\mu$m at
20 fs, 21 $\mu$m at 100 fs), so its depth of focus is shorter than with continuous beams. With $D$ the
continuous-beam depth of focus, the pulsed one is $DG/\sqrt{D^2+G^2}$.
\item The shortening is larger at lower NA, because the focus lengthens while the stretch where the pulses meet
stays the same. At NA 0.5 the depth of focus is 3.56 $\mu$m for continuous beams, 2.81 $\mu$m at 20 fs and
1.92 $\mu$m at 10 fs for the 350-nm pairs (scalar model); with all pairs at 10 fs it is 2.14 $\mu$m through the
fringe crests.
\item The pulse width from which the profiles match the continuous-beam ones to 1\% therefore grows as NA
decreases: about 25 fs at NA 0.9, 87 fs at NA 0.5, 264 fs at NA 0.3 and 610 fs at NA 0.2. For 2PE it is
shorter (62 fs at NA 0.5), because the 350-nm pairs within one pulse need no overlap.
\item With a delay, the pulses meet at $z = c\Delta t/2n$. At low NA this point stays inside the long focus, so
the peak moves along the axis and is kept: at NA 0.1 and 100 fs it is 0.93 of the peak without delay, 14 $\mu$m
from the focus.
\item At NA 0.5 and above the peak stays at the focus and falls with the temporal overlap, to 0.26 at
$\Delta t = \tau$. For all NAs the peak at a delay is
$\exp[-4\ln 2\,z_m^2/(D^2+G^2)]$ with $z_m = c\Delta t/2n$, which gives 0.52 at NA 0.2 against 0.51 stored.
\item Shorter pulses keep more of the peak at $\Delta t = \tau$, since the meeting point moves a shorter
distance: at NA 0.5 it is 0.26 at 100 fs, 0.41 at 20 fs and 0.87 at 5 fs.
\item Below 20 fs the fringes against the delay are shallow along the axis and deepen at high NA (swing 0.03 at
NA 0.5 and 0.27 at NA 0.95, at 5 fs), because the peak follows the 350-nm fringes in space as the delay moves
them.
\end{enumerate}
\begin{center}\includegraphics[width=0.9\textwidth]{delay_facing_vs_na.png}\end{center}

\subsection*{7.6 Medium, pupil, input and pairs}
\begin{enumerate}
\item In water and oil the same results hold at a higher NA, since the focus depends on $\mathrm{NA}/n$. With
facing lenses at NA 0.5, the profiles match the continuous-beam ones from 87 fs in air, 164 fs in water and
216 fs in oil.
\item A longer depth of focus needs a longer pulse to match continuous beams with facing lenses: at NA 0.5,
139 fs for the Gaussian input and 405 fs for the annular pupil with $\varepsilon = 0.9$, against 87 fs for the
circular pupil.
\item The same pupils keep more of the peak under a delay: at NA 0.3, 100 fs and $\Delta t = \tau$, 0.32 for
the circular pupil, 0.38 for the Gaussian input and 0.80 for the annular pupil with $\varepsilon = 0.9$. This
is the use that Cambaliza \& Saloma propose for annular lenses, here for the timing of the pulses.
\item With continuous beams, 2CE needs both beams, which gives its weak background in scattering media. Below
about 15 fs one beam alone also reaches the 350-nm level, through the 375- and 328-nm pairs: at 10 fs their amplitudes
are 18\% and 14\% of that of the 350-nm pairs, and at 20 fs both are below 1\%.
\end{enumerate}

\subsection*{7.7 In short}
The lens angle decides what a pulse changes. With collinear lenses a pulse changes nothing down to 4 fs, with
crossed lenses a few percent below about 20 fs, and with facing lenses it shortens the 2CE focus, more so at
lower NA. A delay costs signal through the temporal overlap at the three angles, except with facing lenses at
low NA, where it moves the 2CE along the axis instead.

\section*{8. Checks and corrections}
\begin{enumerate}
\item The test suite has 1074 tests, which pass. It includes the quoted numbers of the papers, the Linfoot identity
$F = 2Q - C$, and $\nabla\cdot\mathbf{E} = 0$ for the vector field.
\item The pulse and delay runs use tables of the fields. Before the runs the tables agreed with the original code
to $1.6\times10^{-10}$, and 16 samples of the delay runs recomputed with the original code agree to
$1.8\times10^{-10}$.
\item In the review of 2026-11-10 I read the four papers again and compared each equation with the code. One error
was found and corrected: the longitudinal component of the vector field in the two-beam code had a factor $-2i$
in place of $-2$. It changes 2PE with crossed lenses when the polarization lies in the plane of the two beams. It
does not change the 2CE results or the pulse and delay results. The study that holds that case was run again:
208 of its 5712 widths changed, all of them that case, and none of them is on a slide or the dashboard.
\end{enumerate}

\section*{9. Limits and what remains}
\begin{enumerate}
\item The pulses are transform-limited Gaussians, the 350-nm level is narrow compared with the spectra, and the
fringes assume a fixed phase between the two beams.
\item The pulse results use the polarization perpendicular to the plane of the two beams, and crossed lenses
were also run with the polarization in that plane. For real molecules the absorption of two photons with
perpendicular polarizations differs from that of parallel ones by a factor that depends on the molecule, which
the model does not include.
\item 2CE with pulses through a scattering medium has not been computed. It would combine parts 5 and 7.
\item No comparison with an experiment has been made.
\end{enumerate}

\section*{Where things are}
\begin{center}\small\begin{tabular}{ll}\toprule
Review documents, one per part & \texttt{docs/review/01} to \texttt{07}\\
Field of two beams; pulses & \texttt{mcdo/twobeam.py}; \texttt{mcdo/pulsed2ce.py}\\
Replications & \texttt{scripts/replicate\_cambaliza2000.py}, \texttt{scripts/fig\_lim2002\_*.py}\\
Pulse and delay runs & \texttt{scripts/cambaliza2000\_pulses\_*.py}\\
Decks & \texttt{output/08\_two\_color/cambaliza2000/deck/}\\
Dashboard & \texttt{dashboard/cambaliza/}\\\bottomrule\end{tabular}\end{center}
\end{document}
